% !TEX root = ../hphp-project-report.tex
% Generated from HPHP_Documentazione_Completa_EDITABILE.docx.
% The phase title is in English; the source content is preserved in Italian.
\section{Phase A --- Ethnographic Research and Assessment}

\emph{\textbf{FASE A --- Documento 01}}

Team card

\textbf{Nome del gruppo}

Team Affordance

\textbf{Membro del team}

{\def\LTcaptype{none} % do not increment counter
\begin{longtable}[]{@{}
  >{\raggedright\arraybackslash}p{(\linewidth - 2\tabcolsep) * \real{0.3324}}
  >{\raggedright\arraybackslash}p{(\linewidth - 2\tabcolsep) * \real{0.6676}}@{}}
\toprule\noalign{}
\endhead
\bottomrule\noalign{}
\endlastfoot
\textbf{Campo} & \textbf{Valore} \\
Nome completo & Matteo Boscherini \\
Matricola & 0001239423 \\
E-mail & matteo.boscherini@studio.unibo.it \\
Corso di studi & Informatica Magistrale --- Curriculum B \\
\end{longtable}
}

{\def\LTcaptype{none} % do not increment counter
\begin{longtable}[]{@{}
  >{\raggedright\arraybackslash}p{(\linewidth - 2\tabcolsep) * \real{0.3324}}
  >{\raggedright\arraybackslash}p{(\linewidth - 2\tabcolsep) * \real{0.6676}}@{}}
\toprule\noalign{}
\endhead
\bottomrule\noalign{}
\endlastfoot
\textbf{Campo} & \textbf{Valore} \\
Nome completo & Alessandro Campedelli \\
Matricola & 0001242904 \\
E-mail & alessandr.campedell3@studio.unibo.it \\
Corso di studi & Informatica Magistrale --- Curriculum B \\
\end{longtable}
}

{\def\LTcaptype{none} % do not increment counter
\begin{longtable}[]{@{}
  >{\raggedright\arraybackslash}p{(\linewidth - 2\tabcolsep) * \real{0.3324}}
  >{\raggedright\arraybackslash}p{(\linewidth - 2\tabcolsep) * \real{0.6676}}@{}}
\toprule\noalign{}
\endhead
\bottomrule\noalign{}
\endlastfoot
\textbf{Campo} & \textbf{Valore} \\
Nome completo & Francesco Maria Fuligni \\
Matricola & 0001246029 \\
E-mail & francesco.fuligni@studio.unibo.it \\
Corso di studi & Informatica Magistrale --- Curriculum C: Sistemi e
Reti \\
\end{longtable}
}

\emph{\textbf{FASE A --- Documento 02}}

Client card

\textbf{Nome del cliente}

Regione Emilia-Romagna --- Assessorato alle Politiche per la Salute, in
collaborazione con CUP 2000 S.c.p.A., società in-house che gestisce
tecnicamente i sistemi regionali di prenotazione sanitaria (CUPweb, app
ER Salute, Fascicolo Sanitario Elettronico) per conto delle Aziende USL
e Ospedaliere del territorio.

\textbf{Descrizione del cliente}

Ente pubblico regionale che coordina i servizi sanitari di tutta
l\textquotesingle Emilia-Romagna attraverso le Aziende Sanitarie locali
(AUSL Bologna, AUSL Romagna, AUSL Modena, AUSL Parma, AUSL Reggio
Emilia, AUSL Ferrara e le Aziende Ospedaliero-Universitarie). Il cliente
ha già una presenza digitale matura e consolidata: gestisce il Fascicolo
Sanitario Elettronico regionale, il portale di prenotazione CUPweb,
l\textquotesingle app mobile ER Salute, un call center regionale (800
numbers per singola AUSL) e una rete di sportelli fisici e farmacie
convenzionate (FarmaCUP). Non lavoriamo direttamente per la Regione: la
nostra società di consulenza UX è stata incaricata da un system
integrator che opera per conto di CUP 2000 di proporre un redesign
dell\textquotesingle esperienza di prenotazione, da presentare come case
study per un possibile incarico più ampio.

\textbf{Nome del tool}

CUPweb (portale web, www.cupweb.it) e app mobile ER Salute --- sistema
regionale di prenotazione online delle prestazioni sanitarie
specialistiche (visite ed esami) della Regione Emilia-Romagna.

\textbf{Come accedere al tool}

Via browser su www.cupweb.it, tramite SPID, CIE o credenziali FSE/CUPweb
rilasciate agli sportelli; oppure tramite l\textquotesingle app ER
Salute (iOS/Android), che consente di prenotare/disdire visite per sé
stessi e per i familiari associati (minori o persone tutelate).

\textbf{Perché il cliente dovrebbe voler finanziare un nuovo tool?}

Il canale digitale, pur gratuito e disponibile 24/7, resta
sotto-utilizzato dalla fascia più anziana della popolazione, che
continua a intasare call center e sportelli fisici o a delegare
integralmente la prenotazione a un familiare. Questo genera costi
operativi elevati per il servizio sanitario regionale (personale agli
sportelli, call center), un alto tasso di mancate presentazioni
("no-show", oggetto peraltro di una sanzione economica prevista dalla
L.R. 2/2016), e non risponde agli obiettivi di digitalizzazione della
sanità pubblica previsti dal PNRR (Missione 6). Un redesign mirato sulla
fiducia e sull\textquotesingle autonomia degli utenti anziani, che
valorizzi --- invece di scavalcare --- il ruolo del caregiver familiare,
è coerente con questi obiettivi e giustificabile in termini di riduzione
costi.

\textbf{Product Objectives}

Consentire ai pazienti anziani cronici di partecipare più attivamente
alla gestione dei propri appuntamenti sanitari --- anche quando
assistiti da un familiare --- riducendo l\textquotesingle ansia da
interfaccia, aumentando la fiducia nel canale digitale e trasformando il
caregiver da "esecutore" a "facilitatore/allenatore"
dell\textquotesingle autonomia del proprio caro.

\textbf{Business Goals}

Riduzione dei costi operativi di call center e sportelli fisici;
riduzione del tasso di mancata presentazione agli appuntamenti;
incremento della quota di prenotazioni gestite sul canale digitale, in
linea con i target nazionali di digitalizzazione della sanità pubblica
(PNRR Missione 6, componente 2).

\textbf{Brand identity}

Rafforzare la percezione della sanità pubblica regionale come moderna,
efficiente ma allo stesso tempo vicina e inclusiva verso i cittadini più
fragili --- "la sanità digitale che non lascia indietro nessuno" ---
differenziandosi da una digitalizzazione percepita come fredda o ostile
dalla popolazione anziana.

\textbf{Success metrics}

\begin{itemize}
\item
  Aumento della quota di prenotazioni effettuate (anche in co-presenza
  con un familiare) direttamente da utenti over 70
\item
  Riduzione del tempo medio necessario per completare una prenotazione
\item
  Punteggio SUS del nuovo prototipo superiore a 68 (soglia di buona
  usabilità) sul segmento target
\item
  Riduzione del tasso di abbandono durante il flusso di prenotazione
\item
  Riduzione delle chiamate al call center motivate da difficoltà di
  utilizzo del portale/app
\end{itemize}

\emph{\textbf{FASE A --- Documento 03}}

User Segmentation card

\textbf{Dimensioni di segmentazione scelte}

\begin{itemize}
\item
  Demografica: età, presenza/assenza di una rete familiare di supporto
\item
  Comportamentale: frequenza di utilizzo del servizio, grado di delega
  della prenotazione a terzi
\item
  Psicologica: livello di ansia/sfiducia verso gli strumenti digitali,
  motivazione all\textquotesingle autonomia
\item
  Contestuale: dispositivo disponibile, presenza fisica o meno del
  caregiver al momento del bisogno, distanza abitativa
  caregiver--assistito
\end{itemize}

\textbf{Motivazione della scelta, relazione con il tool del cliente}

Il progetto HPHP richiede di considerare due audience distinte ma
interdipendenti: l\textquotesingle utente diretto (che deve usare il
servizio ma non è ancora un utente attivo, ed è scettico/impaurito) e
l\textquotesingle intermediario (che oggi lo sostituisce integralmente).
Per CUPweb/ER Salute, l\textquotesingle utente diretto naturale è il
paziente anziano cronico, che deve prenotare ricorrentemente visite ed
esami; l\textquotesingle intermediario naturale è il caregiver
familiare, che già oggi gestisce gran parte delle pratiche sanitarie del
proprio caro. Le dimensioni scelte permettono di distinguere non solo
per età anagrafica, ma per il reale grado di autonomia/dipendenza
digitale delle due popolazioni, che è il fattore che il redesign deve
realmente affrontare.

\textbf{Segmenti nella dimensione \#1 --- Utenti diretti (per età e
autonomia digitale)}

\begin{itemize}
\item
  1. "Digital newcomer" --- 65-74 anni, competenze digitali di base
  (messaggistica, videochiamate) ma bloccato/a su procedure strutturate
  (SPID, moduli, pagamenti online)
\item
  2. "Digital avoider" --- 75-85+ anni, competenza tecnica quasi nulla,
  delega totale delle pratiche digitali al caregiver, spesso patologie
  croniche con necessità di prenotazioni ricorrenti
\end{itemize}

\textbf{Segmenti nella dimensione \#2 --- Intermediari (per relazione
con l\textquotesingle assistito)}

\begin{itemize}
\item
  1. Caregiver convivente --- vive con l\textquotesingle assistito,
  gestisce le prenotazioni "di persona", spesso presente alle visite
\item
  2. Caregiver "always-on" a distanza --- figlio/a 45-60 anni che non
  convive con il genitore, gestisce le pratiche da remoto (telefono/PC),
  spesso lavoratore/lavoratrice a tempo pieno con poco tempo a
  disposizione
\end{itemize}

\textbf{Numero totale di combinazioni di segmento}

2 (utenti diretti) × 2 (intermediari) = 4 combinazioni possibili.

\textbf{Combinazioni di segmentazione scelte}

\begin{itemize}
\item
  A. Paziente anziano "digital avoider", 75-85 anni, patologia cronica,
  con caregiver figlio/a disponibile ma non sempre fisicamente presente
  (utente diretto)
\item
  B. Caregiver "always-on" a distanza, 45-60 anni, alta competenza
  tecnica, gestisce da remoto le prenotazioni del genitore, vorrebbe
  poter delegare gradualmente compiti semplici (intermediario)
\end{itemize}

\textbf{Perché le combinazioni scelte sono più rappresentative, sfidanti
ed espressive delle altre}

I due segmenti scelti rappresentano i due estremi della relazione di
dipendenza digitale che il progetto HPHP vuole affrontare: da un lato
l\textquotesingle utente completamente dipendente e più a rischio di
esclusione digitale (dato Istat 2025: solo il 63,9\% delle famiglie di
soli anziani ha una connessione internet domestica, contro il 95,1\%
delle altre famiglie), dall\textquotesingle altro
l\textquotesingle intermediario che il cliente vuole trasformare da
"esecutore" a "trainer" dell\textquotesingle autonomia del genitore. È
inoltre il caso più sfidante dal punto di vista del design (bassa
competenza tecnica + alta necessità percepita di
controllo/rassicurazione da parte del caregiver a distanza) e più
rilevante numericamente: in Italia sono oltre 7 milioni le persone che
svolgono attività di caregiving familiare almeno una volta a settimana,
per il 58-60\% donne tra i 45 e i 64 anni.

\textbf{Impatto atteso sul design}

Sarà necessaria un\textquotesingle esperienza a due livelli, con
percorsi guidati passo-passo e a basso carico cognitivo per
l\textquotesingle utente anziano (linguaggio semplice, conferme
esplicite, possibilità di richiamare assistenza umana in ogni momento),
e strumenti di co-prenotazione/delega controllata per il caregiver a
distanza (visibilità sullo stato delle prenotazioni del genitore,
notifiche, possibilità di prenotare per conto di un familiare associato
mantenendo però il genitore informato e coinvolto). Il design dovrà
rendere visibile e graduale il passaggio di competenza dal caregiver al
paziente, coerentemente con l\textquotesingle obiettivo di business del
cliente.

\emph{\textbf{FASE A --- Documento 04 (1/2)}}

Market research card --- Segmento A

\textbf{Segmento di scelta}

Paziente anziano "digital avoider", 75-85 anni, patologia cronica,
utente diretto del servizio, con caregiver familiare disponibile ma non
sempre fisicamente presente.

\textbf{Motivazione della scelta}

È il segmento più esposto al rischio di esclusione digitale
nell\textquotesingle accesso ai servizi sanitari essenziali, e quello su
cui un redesign centrato sulla fiducia e sull\textquotesingle autonomia
deve dimostrare il maggiore impatto.

\textbf{Metodo di ricerca utilizzato --- Available statistics}

Fonti: ISTAT, report "Cittadini e ICT" --- Anno 2024 (pubblicato 28
aprile 2025) e Anno 2025 (pubblicato aprile 2026), istat.it;
approfondimenti giornalistici di sintesi (Sanità Informazione,
AboutPharma) che commentano gli stessi dati Istat.

\begin{itemize}
\item
  Nel 2025 il 63,9\% delle famiglie composte esclusivamente da anziani
  (65+) dispone di connessione a Internet da casa (era il 60,6\% nel
  2024), contro il 95,1\% delle famiglie non esclusivamente anziane e il
  98,7\% delle famiglie con almeno un minore.
\item
  Tra gli ultra 75enni il divario di genere nell\textquotesingle uso di
  Internet è marcato: 38,3\% degli uomini contro il 26,5\% delle donne
  (dato 2024).
\item
  Il 57,8\% delle famiglie senza connessione a Internet indica come
  motivo principale la mancanza di capacità di utilizzo, non il costo
  dei dispositivi o la mancanza di interesse.
\item
  Nel 2024 l\textquotesingle uso di siti/app della Pubblica
  Amministrazione per scaricare moduli è calato di 9,6 punti percentuali
  e per prenotare appuntamenti di 8 punti percentuali rispetto
  all\textquotesingle anno precedente, un segnale di complessità o
  insoddisfazione d\textquotesingle uso più che di disinteresse.
\item
  Il titolo di studio incide fortemente sull\textquotesingle uso di
  Internet: 95,8\% tra i laureati, 91,6\% tra i diplomati, ma solo
  69,7\% tra chi ha al massimo la licenza media --- titolo di studio
  frequente nella coorte 75-85 anni.
\end{itemize}

\emph{(screenshot delle pagine dei report ISTAT da allegare al materiale
supplementare)}

\textbf{DA COMPLETARE --- Interviste one-to-one / osservazione sul campo
con 3-5 pazienti anziani reali (anche solo telefoniche o in presenza di
un familiare), per validare qualitativamente questi dati quantitativi.
Traccia di intervista proposta: (1) Come prenoti di solito una visita
medica? (2) Hai mai provato CUPweb o l\textquotesingle app ER Salute?
Cosa è successo? (3) Chi ti aiuta di solito con queste pratiche? (4)
Cosa ti spaventa di più nell\textquotesingle usare il computer/telefono
per la sanità? (5) Cosa ti farebbe sentire più sicuro/a nel farlo da
solo/a?}

\textbf{Risultati complessivi dell\textquotesingle analisi}

I dati confermano quantitativamente ciò che il progetto HPHP presume
qualitativamente: gli anziani over 75, soprattutto donne e con basso
titolo di studio, sono la fascia più esposta a esclusione digitale nei
servizi essenziali come la sanità. La barriera principale non è
l\textquotesingle assenza di device o di connessione, ma la percezione
di mancanza di competenze --- un problema di fiducia e auto-efficacia
più che di accesso tecnico.

\textbf{Impatto sul design}

Il design deve minimizzare il carico cognitivo e tecnico richiesto
(evitare di esporre da subito procedure complesse come
l\textquotesingle autenticazione SPID), offrire percorsi guidati
passo-passo con conferme esplicite, validare esplicitamente il ruolo del
caregiver come alleato e non sostituto, e prevedere sempre una via di
fuga verso l\textquotesingle assistenza umana (call center, farmacia,
sportello) senza che questa venga percepita come un fallimento.

\emph{\textbf{FASE A --- Documento 04 (2/2)}}

Market research card --- Segmento B

\textbf{Segmento di scelta}

Caregiver familiare "always-on" a distanza, 45-60 anni, alta competenza
tecnica, gestisce da remoto le prenotazioni sanitarie di un genitore
anziano, intermediario del servizio.

\textbf{Motivazione della scelta}

È il segmento che oggi sostiene il costo (tempo, stress, carico mentale)
della mediazione digitale per conto dell\textquotesingle utente diretto,
ed è quello che il cliente vuole trasformare da "esecutore" a
"facilitatore/allenatore" dell\textquotesingle autonomia del proprio
caro --- un obiettivo centrale del progetto HPHP.

\textbf{Metodo di ricerca utilizzato --- Available statistics}

Fonti: rapporto CNEL-Censis "Il valore sociale del caregiver" (per la
Regione Lazio, ottobre 2024); sintesi Istat/EHIS su caregiver familiari
riprese da AboutPharma (2026) e Associazione Carer; dati Istat sulla
popolazione anziana al 1° gennaio 2026.

\begin{itemize}
\item
  In Italia sono oltre 7 milioni le persone che svolgono attività di
  caregiving familiare non retribuito, assistendo un anziano, un
  disabile o un familiare non autosufficiente almeno una volta a
  settimana.
\item
  Il 58-60\% dei caregiver familiari è composto da donne di età compresa
  tra 45 e 64 anni.
\item
  Circa il 40\% dei caregiver, pur essendo in età lavorativa, dichiara
  di non riuscire a conciliare lavoro e attività di assistenza, con
  rischio di impoverimento del nucleo familiare.
\item
  Al 1° gennaio 2026 gli over 65 in Italia sono quasi 14,8 milioni
  (25,1\% della popolazione) e gli over 85 hanno superato i 2,5 milioni,
  con una crescita costante attesa nei prossimi decenni: il fenomeno del
  caregiving è quindi strutturale e non residuale.
\end{itemize}

\emph{(screenshot dei report citati da allegare al materiale
supplementare)}

\textbf{DA COMPLETARE --- Intervista/focus group con 2-3 caregiver
familiari reali (anche colleghi, conoscenti o parenti disponibili,
dichiarandolo onestamente come da istruzioni del corso). Traccia
proposta: (1) Come gestisci oggi le prenotazioni sanitarie di tuo
padre/madre? (2) Quanto tempo/energia richiede questa attività? (3) Hai
mai provato a far prenotare direttamente il tuo familiare? Cosa è
successo? (4) Cosa ti darebbe fiducia per lasciargli più autonomia? (5)
Che tipo di visibilità/controllo vorresti mantenere comunque?}

\textbf{Risultati complessivi dell\textquotesingle analisi}

Il caregiving familiare in Italia è un fenomeno di massa,
prevalentemente femminile e concentrato nella fascia 45-64 anni, con un
costo personale significativo in termini di tempo e possibilità
lavorative. Questo rende particolarmente rilevante, anche dal punto di
vista del cliente (riduzione indiretta dei costi sociali), un design che
alleggerisca il carico del caregiver senza however toglierne il
controllo percepito.

\textbf{Impatto sul design}

Il design dovrà offrire strumenti di co-gestione (visibilità sulle
prenotazioni del familiare, notifiche, possibilità di intervenire solo
quando necessario) piuttosto che di sostituzione totale, con una
progressione esplicita e graduale di autonomia delegata al paziente, in
modo che il caregiver possa fidarsi di "lasciare andare" un compito alla
volta.

\emph{\textbf{FASE A --- Documento 05a}}

Assessment of existing resource card --- Tool del cliente

\textbf{Nome del tool}

CUPweb (portale, www.cupweb.it) e app mobile ER Salute --- Regione
Emilia-Romagna.

\textbf{Piattaforme}

Sito web responsive-limitato (desktop e browser mobile); app nativa
Android e iOS (una versione Windows Phone risulta dismessa).

\textbf{Scopo principale, operazioni consentite, adeguatezza percepita
agli standard professionali}

Il tool consente di prenotare, disdire e modificare appuntamenti per
prestazioni sanitarie specialistiche (SSN e Libera Professione), pagare
online una prenotazione, consultare il Fascicolo Sanitario Elettronico e
ristampare i promemoria degli appuntamenti. Dal punto di vista tecnico
il sito è funzionante e integrato con gli standard nazionali di identità
digitale (SPID, CIE, CNS), ma l\textquotesingle impostazione grafica e
architetturale risulta datata rispetto agli standard di design
contemporanei della Pubblica Amministrazione (confrontabile con le linee
guida di design.gov.it): navigazione a menu testuali densi, assenza di
una gerarchia visiva chiara, servizi distribuiti su domini diversi e
visivamente scollegati tra loro (cupweb.it, pagonlinesanita.it,
fascicolo-sanitario.it, support.fascicolo-sanitario.it).

\textbf{Principali punti di forza, motivo della scelta}

\begin{itemize}
\item
  Gratuito e disponibile 24/7, riduce la necessità di recarsi
  fisicamente a uno sportello
\item
  Integrazione con gli standard nazionali di identità digitale
  (SPID/CIE/CNS), coerente con l\textquotesingle infrastruttura pubblica
\item
  Copertura territoriale capillare su tutta la Regione Emilia-Romagna,
  con canali alternativi complementari (farmacie FarmaCUP, call center,
  sportelli fisici) per chi non può/vuole usare il digitale
\item
  È il vero tool che il cliente ci ha chiesto di ridisegnare:
  analizzarlo nella sua forma reale, difetti inclusi, è indispensabile
  per un redesign credibile
\end{itemize}

\textbf{Principali punti di debolezza}

\begin{itemize}
\item
  Autenticazione obbligatoria complessa (SPID livello 2/CIE) anche solo
  per prenotare, con re-autenticazioni periodiche forzate che gli utenti
  reali segnalano come fonte di frustrazione ricorrente nelle recensioni
  pubbliche dell\textquotesingle app
\item
  Servizi distribuiti su più domini/sottodomini esterni tra loro
  visivamente scollegati, un pattern che in un utente anziano diffidente
  può generare sospetto di phishing invece di fiducia
\item
  Problemi di accessibilità segnalati da utenti con lettori di schermo,
  nonostante la dichiarazione di accessibilità pubblicata
\item
  Interfaccia informativa fortemente testuale, con bassa gerarchia
  visiva e nessun elemento di orientamento per un utente al primo
  utilizzo
\item
  Le funzioni disponibili senza autenticazione sono minime (solo
  pagamento e disdetta), il che costringe da subito
  l\textquotesingle utente meno esperto ad affrontare la parte più
  complessa del flusso (identità digitale) prima ancora di capire il
  valore del servizio
\end{itemize}

First Impressions --- Direct Analysis

\emph{Analisi diretta secondo le categorie delle User Focus Guidelines
(\url{https://www.userfocus.co.uk/resources/guidelines.html}): Navigation, Search, Content
organisation \& layout, Content and titles, Page layout, Shopping/task
completion (qui: prenotazione), Trust and credibility. Il foglio di
valutazione completo con i punteggi guideline-per-guideline va compilato
nel file Excel dedicato; qui riportiamo la sintesi delle criticità più
rilevanti riscontrate.}

\textbf{Elenco delle criticità principali}

\begin{itemize}
\item
  {[}Navigation{]} I link di primo livello ("Disdici", "Prenota",
  "Paga", "Informazioni") portano a domini esterni differenti senza
  preavviso visivo --- viola la linearità di navigazione attesa
\item
  {[}Trust and credibility{]} Il passaggio tra domini diversi (cupweb.it
  → pagonlinesanita.it, fascicolo-sanitario.it) senza un design coerente
  indebolisce la percezione di affidabilità, un fattore critico proprio
  per l\textquotesingle utente anziano diffidente che è il target del
  progetto
\item
  {[}Page layout{]} Assenza di elementi visivi (icone, immagini, colori
  distintivi) che aiutino il riconoscimento rapido delle azioni
  disponibili; il testo è l\textquotesingle unico veicolo informativo
\item
  {[}Content and titles{]} Terminologia tecnico-burocratica non spiegata
  in linguaggio semplice (es. "CNS", "SPID", "impegnativa
  dematerializzata") esposta senza tooltip o spiegazioni contestuali
\item
  {[}Task completion{]} Il completamento della maggior parte delle
  azioni richiede autenticazione forte fin dal primo passo, senza un
  percorso dimostrativo o guidato per il nuovo utente
\item
  {[}Accessibility{]} Segnalazioni pubbliche di incompatibilità con
  screen reader per utenti non vedenti, nonostante
  l\textquotesingle esistenza di una dichiarazione di accessibilità
  formale
\end{itemize}

First Impressions --- Reverse Analysis

\emph{Analisi condotta simulando il task "disdire un appuntamento già
fissato, partendo da un link ricevuto per SMS, senza credenziali SPID
già pronte" --- scenario realistico per il segmento A (paziente
anziano).}

\textbf{Elenco delle criticità principali}

\begin{itemize}
\item
  Il task di disdetta è tra i pochi eseguibili senza autenticazione
  forte, ma non è immediatamente distinguibile in home page da quello di
  prenotazione --- richiede comunque lettura attenta dei menu
\item
  Non è chiaro, senza leggere il testo esplicativo per intero, quale sia
  la differenza fra i canali "Servizio Sanitario Nazionale" e "Libera
  Professione" al momento della scelta, pur essendo un bivio obbligato
  del flusso
\item
  Recensioni reali dell\textquotesingle app ER Salute segnalano che la
  cronologia dei documenti/appuntamenti non è ordinabile o ricercabile
  per data, costringendo a scorrimenti lunghi --- un problema di
  efficienza soprattutto per un utente con minore dimestichezza con lo
  scroll touch
\item
  Il cambio dell\textquotesingle obbligo di autenticazione (introduzione
  dell\textquotesingle obbligo SPID dove prima bastava il riconoscimento
  biometrico) è stato vissuto da utenti reali come un peggioramento
  improvviso e non comunicato in modo comprensibile, minando la fiducia
  già costruita
\end{itemize}

\emph{\textbf{FASE A --- Documento 05b}}

Assessment of existing resource card --- Tool competitor

\textbf{Nome del tool}

MioDottore (Doctolib Italia) --- piattaforma privata di prenotazione di
visite mediche.

\textbf{Piattaforme}

Sito web (miodottore.it) responsive, app nativa Android e iOS.

\textbf{Scopo principale, operazioni consentite, adeguatezza percepita
agli standard professionali}

Piattaforma privata (non SSN) che permette di cercare oltre 240.000
medici e specialisti per categoria, città o assicurazione sanitaria,
leggere recensioni di altri pazienti, prenotare una visita in pochi
passaggi, ricevere promemoria, gestire/annullare appuntamenti e
contattare lo specialista via messaggi o videoconsulto. Il design è
moderno, coerente con gli standard commerciali attuali (ricerca a
faccette, mappa, card con foto, iconografia chiara).

\textbf{Principali punti di forza, motivo della scelta}

\begin{itemize}
\item
  Un solo dominio, un solo brand coerente per tutto il flusso --- nessun
  salto percepito di fiducia come invece accade in CUPweb
\item
  Registrazione/accesso semplificato (email o numero di telefono), senza
  richiedere un\textquotesingle identità digitale certificata come SPID
  per la sola consultazione o prenotazione
\item
  Interfaccia fortemente visiva: foto dei professionisti, valutazioni a
  stelle, filtri chiari, mappa --- riduce il carico cognitivo rispetto
  al solo testo
\item
  Promemoria automatici e possibilità di messaggistica diretta col
  professionista, elementi di rassicurazione assenti in CUPweb
\item
  È stato scelto come termine di paragone perché, pur avendo uno scopo
  diverso (mercato privato vs. sanità pubblica), rappresenta un
  benchmark di riferimento per gli standard di semplicità che un utente
  anziano o il suo caregiver oggi si aspettano da un\textquotesingle app
  di prenotazione sanitaria, avendola probabilmente già vista usare da
  altri familiari
\end{itemize}

\textbf{Principali punti di debolezza}

\begin{itemize}
\item
  Le recensioni sono moderate secondo linee guida della piattaforma
  stessa, il che ha generato dubbi pubblici sulla loro
  genuinità/affidabilità --- un problema di trasparenza che il progetto
  HPHP (dove la fiducia è centrale) deve evitare di replicare
\item
  Alcuni utenti segnalano un flusso di login inutilmente frammentato tra
  browser, modali e rimando all\textquotesingle app store, con
  conseguente frustrazione
\item
  Il servizio è privato: non integra il Servizio Sanitario Nazionale, le
  esenzioni ticket, o il Fascicolo Sanitario Elettronico --- quindi non
  risolve da solo il caso d\textquotesingle uso del cliente, ma resta
  valido come fonte di ispirazione sul piano
  dell\textquotesingle esperienza
\item
  Il modello di business (i medici pagano per essere visibili) può
  introdurre bias nei risultati di ricerca, un problema di trasparenza
  inapplicabile ma da tenere presente come contro-esempio
\end{itemize}

First Impressions --- Direct Analysis

\textbf{Elenco delle criticità principali}

\begin{itemize}
\item
  {[}Trust and credibility{]} Il sistema di recensioni moderate genera
  sospetto in una parte degli utenti, minando proprio
  l\textquotesingle elemento di fiducia che dovrebbe essere il punto di
  forza
\item
  {[}Task completion{]} Il login, secondo segnalazioni reali, può
  risultare inutilmente lungo (redirect multipli tra app e browser) su
  alcuni dispositivi
\item
  {[}Navigation{]} Buona nel complesso: ricerca a faccette e mappa
  riducono il carico di navigazione libera, un pattern che vale la pena
  studiare per il redesign
\end{itemize}

First Impressions --- Reverse Analysis

\emph{Analisi condotta simulando il task "trovare e prenotare una visita
specialistica vicino a casa, come utente al primo accesso".}

\textbf{Elenco delle criticità principali}

\begin{itemize}
\item
  Il primo accesso è comunque mediato da un passaggio di registrazione,
  anche se più leggero di SPID --- un piccolo attrito resta anche nel
  miglior competitor osservato
\item
  La grande quantità di filtri e opzioni (assicurazione, videoconsulto,
  tipo di visita) può risultare eccessiva per un utente che cerca
  semplicemente "la prossima visita disponibile", un rischio di
  sovraccarico da tenere presente anche nel nostro redesign
\end{itemize}

\emph{\textbf{FASE A --- Documento 06a}}

Preliminary User testing card --- Tool del cliente

\emph{Questo test riguarda CUPweb / app ER Salute, il tool da
ridisegnare.}

\textbf{Nome del tool}

CUPweb (portale) / app ER Salute

Protocollo di test

\textbf{Setup della sessione di test}

Test in presenza, presso il domicilio del soggetto o in un luogo neutro
familiare al soggetto (per ridurre l\textquotesingle ansia da contesto),
su dispositivo del soggetto stesso quando possibile (smartphone o PC
abitualmente usato), altrimenti su un dispositivo del team
preconfigurato in modo identico. Un membro del team fa da
assistente/osservatore, un secondo prende appunti.

\textbf{Metodo di test}

\begin{itemize}
\item
  Discount usability testing (3 soggetti, in linea con le indicazioni
  del corso per test preliminari)
\end{itemize}

\textbf{Elenco dei task da testare}

\begin{itemize}
\item
  T1 --- Accedere al portale/app e individuare come prenotare una visita
  specialistica (senza completare la prenotazione)
\item
  T2 --- Disdire un appuntamento fittizio già presente
  nell\textquotesingle account demo/test
\item
  T3 --- Individuare dove e come chiedere aiuto/assistenza in caso di
  difficoltà
\end{itemize}

\textbf{Metodologia di test}

\begin{itemize}
\item
  Thinking aloud per i task T1 e T3 (per comprendere il ragionamento e
  le difficoltà percepite)
\item
  Bottom line data per il task T2 (per misurare efficienza/efficacia
  reale senza interferenze)
\end{itemize}

\textbf{Metriche di test}

\begin{itemize}
\item
  Success (task completato sì/no)
\item
  Efficiency (tempo impiegato)
\item
  Efficacy (accuratezza del risultato ottenuto)
\item
  Errors (numero e tipo di errori commessi)
\item
  Learnability (miglioramento fra primo e secondo tentativo, se
  applicabile)
\item
  Satisfaction (tramite questionario SUS)
\end{itemize}

\textbf{Reclutamento dei soggetti}

\textbf{DA COMPLETARE --- Da completare con soggetti reali. Requisiti
minimi: almeno 2-3 persone che rientrino nel segmento A (65-85 anni,
bassa/media competenza digitale) e/o segmento B (caregiver familiare
45-60 anni). Per ciascun soggetto indicare: nome fittizio, età, genere,
competenza tecnica e di dominio, somiglianza con il segmento scelto,
relazione con il team (es. "nonna di un membro del team", "vicino di
casa"). Almeno due di questi soggetti dovranno essere ricontattati nella
Fase D per valutare il redesign: pianificare fin da ora una seconda
sessione con loro.}

\textbf{Modalità di reclutamento dei soggetti}

\textbf{DA COMPLETARE --- Da completare: descrivere onestamente come
sono stati trovati i soggetti (es. rete familiare/amicale dei membri del
team, associazioni per anziani del quartiere, centri sociali,
parrocchie, ecc.).}

\textbf{Motivazione della scelta dei soggetti}

\textbf{DA COMPLETARE --- Da completare con onestà, come richiesto
esplicitamente dalla card originale: se i soggetti sono stati scelti per
disponibilità/vicinanza più che per rappresentatività statistica, va
dichiarato esplicitamente.}

\emph{Le sezioni "Test \#1: Subject\ldots" (nome fittizio, data,
assistente, contesto, dati qualitativi e quantitativi, punteggio SUS dal
modulo Google fornito dal corso, suggerimenti dei soggetti, riflessioni
utili per il design) andranno aggiunte --- una per ciascun soggetto
reale testato --- non appena i test saranno stati condotti. Il
protocollo sopra è pronto per essere eseguito.}

\emph{\textbf{FASE A --- Documento 06b}}

Preliminary User testing card --- Tool competitor

\emph{Questo test riguarda MioDottore.}

\textbf{Nome del tool}

MioDottore (Doctolib Italia)

Protocollo di test

\textbf{Setup della sessione di test}

Stesso setup del test sul tool del cliente (in presenza, dispositivo
abituale del soggetto quando possibile), idealmente con gli stessi
soggetti reclutati per il test su CUPweb, per poter confrontare le
esperienze a parità di persona.

\textbf{Metodo di test}

\begin{itemize}
\item
  Discount usability testing (3 soggetti)
\end{itemize}

\textbf{Elenco dei task da testare}

\begin{itemize}
\item
  T1 --- Cercare uno specialista di una categoria a scelta nella propria
  città (senza completare la prenotazione)
\item
  T2 --- Individuare come si annullerebbe un appuntamento già fissato
\item
  T3 --- Individuare dove e come si contatta il professionista in caso
  di dubbi
\end{itemize}

\textbf{Metodologia di test}

\begin{itemize}
\item
  Thinking aloud per T1 e T3
\item
  Bottom line data per T2
\end{itemize}

\textbf{Metriche di test}

\begin{itemize}
\item
  Success, Efficiency, Efficacy, Errors, Learnability, Satisfaction
  (SUS)
\end{itemize}

\textbf{Reclutamento dei soggetti}

\textbf{DA COMPLETARE --- Idealmente gli stessi soggetti del test 06a,
per un confronto diretto delle due esperienze. Vedi istruzioni nella
card 06a.}

\textbf{Modalità di reclutamento dei soggetti}

\textbf{DA COMPLETARE --- Vedi card 06a.}

\textbf{Motivazione della scelta dei soggetti}

\textbf{DA COMPLETARE --- Vedi card 06a.}

\emph{Anche qui, le sezioni "Test \#1: Subject\ldots" per ciascun
soggetto reale vanno aggiunte a valle dell\textquotesingle esecuzione
dei test. Protocollo pronto per essere eseguito.}

\emph{\textbf{FASE A --- Documento 07}}

Conclusions and design recommendation card

\textbf{Elenco delle criticità per importanza}

\emph{Elenco preliminare basato sulle sole Assessment card (05a/05b).
Andrà integrato e ri-prioritizzato con impatto/persistenza reali una
volta raccolti i dati dello User testing (06a/06b).}

\begin{itemize}
\item
  {[}ID-01{]} Autenticazione obbligatoria complessa (SPID/CIE) richiesta
  fin dal primo passo, anche per compiti semplici --- impatto: alto;
  persistenza: alta
\item
  {[}ID-02{]} Frammentazione tra domini/sottodomini scollegati tra loro
  (cupweb.it, pagonlinesanita.it, fascicolo-sanitario.it) --- impatto:
  alto (mina la fiducia); persistenza: alta
\item
  {[}ID-03{]} Interfaccia poco gerarchizzata, prevalentemente testuale,
  priva di elementi visivi di orientamento --- impatto: medio;
  persistenza: alta
\item
  {[}ID-04{]} Terminologia tecnico-burocratica non spiegata (SPID, CNS,
  impegnativa dematerializzata) --- impatto: medio; persistenza: media
\item
  {[}ID-05{]} Assenza di un percorso guidato/dimostrativo per il nuovo
  utente --- impatto: alto; persistenza: alta
\item
  {[}ID-06{]} Problemi di accessibilità per utenti con tecnologie
  assistive --- impatto: alto (per il sottoinsieme di utenti coinvolti);
  persistenza: media
\item
  {[}ID-07{]} Cronologia appuntamenti/documenti non ordinabile o
  ricercabile nell\textquotesingle app --- impatto: basso-medio;
  persistenza: media
\end{itemize}

\textbf{Tabella riassuntiva dei punteggi SUS}

\textbf{DA COMPLETARE --- Da compilare con i punteggi reali ottenuti dal
modulo Google SUS somministrato ai soggetti di test (card 06a e 06b).
Struttura pronta: una colonna per soggetto, righe per le 10 domande
standard SUS, riga Totale, Media, Varianza. Se la varianza tra soggetti
risultasse elevata, andranno indagate le ragioni (es. soggetti con
competenze digitali molto diverse tra loro) e valutata
l\textquotesingle opportunità di soggetti aggiuntivi.}

\textbf{Curva di urgenza}

\textbf{DA COMPLETARE --- Da costruire (almeno impatto × persistenza)
una volta confermate le criticità con i dati di test reali, usando gli
identificatori {[}ID-01...ID-07{]} sopra elencati come base di partenza,
integrati con quanto emerso dai test.}

\textbf{Raccomandazioni di design preliminari}

\emph{Raccomandazioni preliminari, connesse alle criticità individuate
nell\textquotesingle assessment. Andranno arricchite/confermate con le
evidenze dello user testing prima della stesura finale.}

\begin{itemize}
\item
  {[}R-01{]} (collegata a ID-01, ID-05) Introdurre un percorso "a bassa
  soglia" iniziale che permetta di esplorare il servizio (es. vedere
  disponibilità, capire come funziona) prima di richiedere
  l\textquotesingle autenticazione forte, spostando SPID/CIE solo al
  momento realmente necessario
\item
  {[}R-02{]} (collegata a ID-02) Unificare visivamente tutti i
  touchpoint del servizio (prenotazione, pagamento, fascicolo) sotto
  un\textquotesingle unica identità grafica coerente, anche se
  tecnicamente restano sistemi diversi, per non rompere la fiducia
  dell\textquotesingle utente durante il percorso
\item
  {[}R-03{]} (collegata a ID-03, ID-04) Introdurre gerarchia visiva
  (icone, colori, dimensioni) e microcopy in linguaggio semplice con
  spiegazioni contestuali per i termini tecnici (tooltip, esempi)
\item
  {[}R-04{]} (collegata a ID-05) Progettare un flusso guidato
  passo-passo per il primo utilizzo, con conferme esplicite ad ogni
  passaggio e possibilità di richiamare assistenza (caregiver o
  operatore) senza percepirlo come un fallimento
\item
  {[}R-05{]} (collegata a ID-06) Verificare e correggere la
  compatibilità con tecnologie assistive (screen reader) in coerenza con
  la dichiarazione di accessibilità già pubblicata dal cliente
\item
  {[}R-06{]} (collegata al segmento B / caregiver) Introdurre una
  modalità di "co-gestione" che dia visibilità al caregiver sulle
  prenotazioni del familiare pur lasciando che sia
  quest\textquotesingle ultimo, quando possibile, a compiere
  l\textquotesingle azione finale --- per costruire gradualmente
  autonomia invece di sostituirla
\end{itemize}
